\section{Introduction}\label{sec:introduction}

Hilbert's tenth problem over the rational numbers asks for a decision
procedure for rational zeros of polynomials with integer coefficients.
The number of variables is part of the input.

\begin{theorem}\label{thm:main}
There is no algorithm which, given a polynomial
$f\in\Z[X_1,\ldots,X_n]$, with $n$ part of the input, decides whether
$f$ has a zero in $\Q^n$.
\end{theorem}

\subsection{Historical setting}

Hilbert's tenth problem, as stated in his 1900 list of mathematical
problems, asks for a finite procedure deciding whether a polynomial
equation with integer coefficients has a solution in the ordinary
integers \cite[Problem~10]{Hilbert1900}.  Its negative resolution grew
out of a much stronger description of polynomial equations.  A relation
on the positive integers is \emph{recursively enumerable} if an
algorithm can list exactly its satisfying tuples.  It is
\emph{Diophantine} if membership can be expressed by the existence of
additional positive integers satisfying a polynomial equation with
integer coefficients.  Thus a Diophantine description expresses a
condition entirely through polynomial solvability.

Davis, Putnam, and Robinson proved in 1961 that every recursively
enumerable relation has such a representation if exponentiation is
also allowed.  They reduced the ordinary polynomial representation
theorem to a growth criterion for a Diophantine relation
\cite[Theorem and Corollary~3]{DavisPutnamRobinson1961}.
Matiyasevich supplied the required growth in 1970 through a
Diophantine description of a Fibonacci relation, completing the
theorem that every recursively enumerable relation on the positive
integers is Diophantine \cite{Matiyasevich1970}.  The equivalent
all-integer decision problem is therefore undecidable.  This
representation theorem explains the strength of the result: integer
polynomial equations can encode arbitrary recursively enumerable
conditions.

Letting the unknowns range over \(\mathbb Q\) changes the decision
problem.  Enumerating rational tuples will eventually find any rational
zero, but gives no answer for an equation without one.  A possible
transfer of the integer theorem would be an existential definition of
membership in \(\mathbb Z\) using only the ring operations of
\(\mathbb Q\).  One could then replace each integer unknown by a
rational unknown subject to that membership condition, obtaining an
equivalent rational existential question.  More generally, a
\emph{Diophantine model} of integer arithmetic represents the integers
by a Diophantine subset of some \(\mathbb Q^k\), with Diophantine
relations for addition and multiplication.  Such a model also
translates integer equations into rational existential questions
\cite[Definition~3.1 and Observation~3.3]{CornelissenZahidi2000}.
Here the defining polynomial equations and all their witnesses are
interpreted over \(\mathbb Q\).
These constructions are sufficient routes to undecidability; the
decision problem does not require its solution to produce either one.

Julia Robinson showed in 1949 that the integers are first-order
definable in the rationals using addition and multiplication.  Her
formula uses both universal and existential quantifiers, and her work
also proves undecidability of the full first-order theory of rational
arithmetic \cite[Theorems~3.1 and~4.2]{Robinson1949}.  Later work
reduced the quantifier complexity.  Poonen defined \(\mathbb Z\) in
\(\mathbb Q\) by a positive formula with two universal quantifiers
followed by seven existential quantifiers
\cite[Theorem~4.1]{Poonen2009}.  Koenigsmann then gave a universal
definition of \(\mathbb Z\); equivalently, its complement in
\(\mathbb Q\) is Diophantine.  His result also implies undecidability
of the \(\forall\exists\) theory
\cite[Theorem~1 and Corollaries~2--3]{Koenigsmann2016}.  These results
give substantial control of definitions inside \(\mathbb Q\), but
do not supply the existential membership condition for the direct
transfer above.

The distinction between a Diophantine model and undecidability also
has a geometric setting.  Mazur conjectured that, for a variety over
\(\mathbb Q\), the closure of its rational points in its real points
has only finitely many connected components
\cite[Conjecture~3]{Mazur1995}.  He observed that this conjecture
would rule out a Diophantine definition of \(\mathbb Z\) in
\(\mathbb Q\).  Cornelissen and Zahidi showed that it would also rule
out a Diophantine model of integer arithmetic in a finite Cartesian
power of \(\mathbb Q\)
\cite[Theorem~3.6]{CornelissenZahidi2000}.  These obstructions to
definitions and models do not by themselves rule out a negative
answer to the rational decision problem.

One approach studies the subrings
\(\mathbb Z[S^{-1}]\subseteq\mathbb Q\), whose denominators use only
primes in \(S\).  Poonen constructed disjoint recursive sets of primes
\(T_1,T_2\), each of natural density zero, such that positive-integer
arithmetic has a Diophantine model in \(\mathbb Z[S^{-1}]\) whenever
\(S\supseteq T_1\) and \(S\cap T_2=\varnothing\).  Hilbert's tenth
problem is undecidable for every such ring.  Taking \(S\) to be the
complement of \(T_2\) gives a recursive set of density one and a
proper subring of \(\mathbb Q\)
\cite[Theorem~1.3]{PoonenSubrings}.  The construction uses selected
multiples of a point on a rank-one elliptic curve whose real
coordinates approximate integers.  

Eisenträger, Miller, Park, and Shlapentokh later constructed
computably presentable subrings \(\mathbb Z[S^{-1}]\), including
examples in which \(S\) has density one, whose Hilbert's tenth problem
is Turing equivalent to that over \(\mathbb Q\)
\cite[Theorem~3.17]{EisentragerMillerParkShlapentokh2017}.  For these
examples, a Turing reduction from integer solvability would already
give a Turing reduction from integer solvability to the rational
field problem.  The density of the
inverted primes alone therefore does not settle the relation to the
field problem.

Elliptic curves also give a transfer between rings of integers of
number fields.  Poonen proved that, for \(F\subseteq K\), an elliptic
curve \(E/F\) with
\[
             \operatorname{rank}E(F)=\operatorname{rank}E(K)=1
\]
makes \(\mathcal O_F\) Diophantine in \(\mathcal O_K\)
\cite[Theorem~1]{PoonenRankOne2002}.  If \(\mathbb Z\) is already
Diophantine in \(\mathcal O_F\), undecidability transfers to
\(\mathcal O_K\).  Mazur and Rubin pursued this route through
quadratic twists and Selmer groups, conditional on the conjectured
evenness of \(\dim_{\mathbb F_2}\operatorname{Sha}(E/L)[2]\).
Assuming this for every elliptic curve over every number field, they obtained undecidability for
every infinite finitely generated \(\mathbb Z\)-algebra
\cite[Theorem~1.13, Corollary~1.14, and Theorem~8.1]{MazurRubin2010}.

This number-field program has recently been completed
unconditionally.  Koymans and Pagano combined two-descent with
additive combinatorics to prove that \(\mathbb Z\) is Diophantine in
the ring of integers of every number field, and that Hilbert's tenth
problem is undecidable for every infinite finitely generated
\(\mathbb Z\)-algebra \cite{KoymansPagano2026}.
Alpöge, Bhargava, Ho, and Shnidman
gave another route: for every quadratic extension \(K/F\) of number fields, they
construct an abelian variety over \(F\) with positive rank unchanged
over \(K\), and deduce that \(\mathbb Z\) is Diophantine in every
number-field ring of integers
\cite[Theorem~1.1 and Corollary~1.2]{AlpogeBhargavaHoShnidman2026}.
The field \(\mathbb Q\) is not a finitely generated
\(\mathbb Z\)-algebra, so these results concern a different class of
rings from the field in Theorem~\ref{thm:main}.

The elliptic approach also illuminates nearby rational decision
problems.  In a rank-one group, a multiple \(nP\) of a nontorsion
point can represent the integer \(n\), and local coordinates or heights can reveal arithmetic
information about that index.  Cornelissen and Zahidi used elliptic
divisibility sequences and existential valuation predicates in this
direction.  Under a conjectural odd primitive-divisor property for one
such sequence, they obtained a Diophantine model of
\((\mathbb Z,+,\mid)\) and proved undecidability of a positive
\(\Pi_2\) fragment with one universal quantifier
\cite[Conjecture~3.16 and Theorems~3.17 and~5.3]{CornelissenZahidi2007}.
The assertion with one universal quantifier uses the correction by
Cornelissen and Shlapentokh
\cite[Remark~2.3 in the author version]{CornelissenShlapentokh2009}.


Garcia-Fritz, Pasten, and Vidaux instead adjoined predicates comparing
the logarithmic heights of rational tuples to the rational ring
language.  In this expanded language they gave a
positive existential interpretation of \(\mathbb N\), and hence
proved undecidability; comparisons between tuples of length at most
three suffice \cite[Theorems~1.1--1.2]{GarciaFritzPastenVidaux2025}.
Their construction uses elliptic canonical heights to encode
consecutive squares.  The present paper concerns ordinary rational
polynomial solvability, and its height estimates enter the proof of
that ordinary statement.

\subsection{Proof strategy}

Constant polynomials can be decided directly. For each nonconstant integral
polynomial \(f\), we construct an effectively generated sequence of finite tests.
Each test can be answered by finitely many rational-solvability queries, and the sequence satisfies
\begin{equation}\label{intro:finite-test-interface}
 f\text{ has an integer zero}
 \quad\Longleftrightarrow\quad
 \text{every test succeeds}.
\end{equation}
An integer zero supplies witnesses for all the tests. If a rational-solvability algorithm existed, we could run two searches in parallel: enumeration of integer zeros, and enumeration of tests until one fails. The arithmetic work is to prove that a failed test exists whenever \(f\) has no integer zero.

Suppose all the finite tests succeed. Compactness then gives a ring \(R\) containing a zero \(\mathbf a=(a_1,\ldots,a_n)\) of \(f\), embedded in an elementary extension of the rational field; we use a tilde for embedded values and points. The ambient extension also carries the integer and valuation data, height functions, elliptic multiple maps, and functions for finite products used in the proof. The tests attach an elliptic point \(T_a\) to each \(a\in R\), compatibly with addition, and require \(T_a\) to be nonidentity when \(a\ne0\). The embedded points lie in the transferred copy of a fixed infinite cyclic subgroup \(\Z P\) of finite index in a rank-one elliptic group, with \(T_1=P\). This defines unique indices
\[
                 \widetilde T_a=\eta(a)P,
                 \qquad \eta(a)\in{}^*\Z.
\]
Here \({}^*\Z\) is the ambient integer structure, whose elements may be nonordinary. The point axioms make \(\eta:R\to{}^*\Z\) injective and additive, with \(\eta(1)=1\), but \(\eta\) need not respect multiplication. Thus the modeled equation does not yet give an equation satisfied by the indices.

Two local comparisons address this obstruction. The primary test compares an embedded ring value with a quotient of truncated elliptic logarithms. The formal logarithm can recover the corresponding quotient of point indices once that quotient is integral at the prime. A second comparison, on the conjugate elliptic curve, supplies this control using only additivity of \(\eta\). Where the comparisons apply, they show that every \(\widetilde a\) is integral at the prime and agrees with \(\eta(a)\) to the prescribed precision. Each pair uses auxiliary choices that work for every ring element, which is essential for the later uniform bound.

The proof arranges these comparisons at primes detected by a height estimate. Fix the five rational points of \(\mathbf P^1\) supplied by that estimate and primitive integral linear forms vanishing at them; call these the contact forms. For a rational point \(s=(u:v)\) outside this set, written in coprime integer coordinates, let \(M(s)\) be the product, without repetition, of the nonexceptional primes at which one of these forms has odd valuation. We call the primes counted by \(M(s)\) its contact primes. The estimate bounds the absolute logarithmic Weil height by
\[
                              h(s)\le H M(s)^c
\]
for fixed constants \(H,c\). This definition and estimate transfer to the ambient extension. At each prescribed precision, the theory represents every ring element as a
sum of three fractions of ring elements, called slopes, and equips factors of
their contact forms with the local conditions. A prime-pattern theorem supplies
the required primes for these representations in the ordinary integer model.
A separate existential condition controls the parity of possible denominator
valuations outside a fixed finite set. In a model satisfying the tests, consider
the three slopes supplied for any ring element at that precision. After the
further fixed exclusions, the parity condition lets the proof apply the local
comparisons at every contact prime of those slopes. Enlarging the exceptional
set changes the height constant but preserves the exponent \(c\).

Set
\[
 \delta=f\bigl(\eta(a_1),\ldots,\eta(a_n)\bigr),
 \qquad
 B=\max\bigl(1,|\eta(a_1)|,\ldots,|\eta(a_n)|\bigr).
\]
If \(\delta=0\), elementarity gives an ordinary integer zero of \(f\). Otherwise \(\delta\) is a nonzero ambient integer whose size is bounded by a fixed polynomial in \(B\). Choose one ordinary precision \(K\) large enough in terms of the degree of \(f\) and the exponent \(c\). For any \(A\in R\), choose a tested representation at precision \(K\). Every contact prime of its slopes divides this same \(\delta\) to order at least \(K\), because all modeled root coordinates agree locally with their indices. The height estimate then bounds each slope height by a constant times \(B\). The logarithmic height inequality for their sum gives
\[
                         h(\widetilde A)\ll_f B
                         \qquad\text{for every }A\in R,
\]
with one ordinary constant.

For each nonzero root coordinate \(a_j\), consider its attached point
\(T_{a_j}\). Applying the uniform bound to the ring elements representing
\(x(T_{a_j})\) gives \(h(x(\eta(a_j)P))\ll_f B\).
The transferred canonical-height comparison for \(P\) makes this height grow
quadratically in \(|\eta(a_j)|\). If \(B>1\), choose \(j\) with
\(|\eta(a_j)|=B\); the two bounds give \(B^2\ll_f B\). If \(B=1\), the
required bound is immediate. Thus all root indices lie in an ordinary finite
interval. Additivity and injectivity identify the modeled root coordinates
with those ordinary integers, contradicting \(\delta\ne0\). This proves the
finite-test equivalence and hence the undecidability theorem.

The detailed reduction in Section~\ref{sec:reduction} makes these steps
precise. Section~\ref{par:section} proves the pole-parity condition; its
passage from two-descent data to rational points uses the pointwise
$2$-converse \cite[Theorem~1.1]{PointwiseTwoConverse} and the
Cassels--Tate pairing \cite{Cassels1962}.
Section~\ref{hei:section} proves the five-point height estimate. Its
modularity step uses the odd regular Fontaine--Mazur theorem at $2$
\cite[Theorem~1.1]{OddRegularFontaineMazurAtTwo}, whose statement has no
hypothesis on the residual image. Section~\ref{ell:section} proves the
rank-one and local elliptic results, and Section~\ref{pri:section}
supplies the prime patterns and simultaneous local witnesses for the
ordinary integer model. Section~\ref{red:degree-normal-form} gives
Turing degree $0'$ and the quartic normal form.

At rational primes, valuations are normalized by \(v_q(q)=1\), with
\(v_q(0)=+\infty\). The support of a nonzero rational number is the set of
primes dividing its numerator or denominator in lowest terms. The notation
\(X\ll Y\) means \(X\le CY\) for a positive constant depending only on the
fixed data in the statement; additional dependencies are indicated when
they occur.
